% Principal GDZ PPN236018647, native471lower–474; HKUST472–475.
\NumberedItem{5}
\begin{center}\textsc{Gauss} an \textsc{Schumacher}.\quad Göttingen, 13. April 1836.\end{center}
\HistoricCitation{[Briefwechsel zwischen \textsc{Gauss} und \textsc{Schumacher} III, Altona 1861, S. 30—32.]}
\indent Wegen der geometrischen Aufgabe bin ich sehr gern erbötig, mich Ihnen näher zu expliciren, sobald Sie, mein theuerster Freund, mir erst näher bezeichnen, was denn eigentlich Ihnen dabei einer nähern Explication bedürftig scheint. Obgleich ich Ihren Brief wiederholt gelesen habe, ist es mir doch nicht möglich gewesen, dieses daraus abzunehmen. Da nun aber der Gegenstand \textls[120]{viele Seiten} hat, und eine \textls[120]{complete} Auseinandersetzung von \textls[120]{Allem} ein kleines Buch füllen könnte, so würde ich, wenn ich auf gut Glück Eine Seite hervorhöbe, und alles was ich darüber sagen könnte, entwickelte, riskiren, mich ganz vergeblich abgequält zu haben, wenn sich nachher fände, dass das gar nicht der Ihnen dunkele Punkt gewesen sei. Sie leiten aus Ihrem
\clearpage\SourceLeaf{472}\ArithmeticHead{465}{R}{GAUSS AN SCHUMACHER.}
\noindent Calcul zwei Gleichungen ab:
\[
\text{(A)}\qquad\left\{\begin{aligned}
0&=\cos w+\cos w'+\cos w''+\cos w'''\ .\ .\ .\\
0&=\sin w+\sin w'+\sin w''+\sin w'''\ .\ .\ .\,.
\end{aligned}\right.
\]
\noindent Sie sagen ferner: diesen Gleichungen geschehe Genüge, wenn die von \((x,y)\) nach \((a,b),\ (a',b'),\ (a'',b'')\ .\ .\ .\) gezogenen Linien lauter gleiche Winkel bilden und für ein Viereck, dessen innere Winkel alle \(<180\) sind, durch den Durchschnitt der Diagonalen. (Ich verstehe dies so, dass Sie zwei Fälle als von selbst einleuchtend vorfinden 1) allgemein wenn \(w'-w=w''-w'=w'''-w''=\text{etc.}=w-w^n\), 2) wenn bei 4 Pkten \(w''=w+180^{0},\ w'''=w'+180^{0}\)). Mehr (fahren Sie fort) kann ich aber nicht aus den Grundgleichungen herausbringen.

\indent Ich weiss nun eigentlich nicht, was Sie denn mehr gewünscht haben. Ich kann also nur rathen. Z. B.

\indent 1) wollen Sie sich nicht mit der Aufgabe, wo nur von 4 Pkten die Rede ist (von welcher allein in unsern frühern Briefen die Rede gewesen ist), beschäftigen, sondern auch mit einer grössern Zahl? Eine solche \textls[120]{allgemein} betrachtet wird immer auf höhere Gleichungen führen müssen,

\indent 2) oder wenn Sie, wie der Anfang Ihres Briefes vermuthen lässt, sich \textls[120]{nur} mit 4 Pkten beschäftigen, was wünschen Sie in diesem Fall mehr? Ist es Ihnen vielleicht nicht einleuchtend, dass dann aus den Gleichungen (A) Ihr Schluss \textls[120]{mit Nothwendigkeit} folgt? Dies würde sich (sobald ich weiss, dass dies Ihr Bedenken ist) leicht erledigen können, oder

\indent 3) finden Sie Anstoss daran, dass das Resultat (Schnitt von 2 Diagonalen) nur für Vierecke gilt, deren Winkel alle \(<180^{0}\), und dass bei vier Punkten wie [die Figur sie zeigt], die ein Dreieck, mit dem vierten Pkte inwendig, bilden, nothwendig dieser vierte Pkt die Auflösung gibt? Dann würde sich die Metaphysik davon und \textls[120]{warum} hier die blinde Rechnung aus den Gleichungen (A) nicht zureicht, leicht aufhellen lassen\quad Oder verlangen Sie
\begin{center}
\begin{tikzpicture}[x=.016cm,y=-.016cm,line width=.45pt]
\draw (77,0) circle[radius=3.3];
\draw (78,56) circle[radius=3.3];
\draw (0,98) circle[radius=3.3];
\draw (149,137) circle[radius=3.3];
\end{tikzpicture}
\end{center}
\par\vfill\noindent X1.\hfill59
\clearpage\SourceLeaf{473}\ArithmeticHead{466}{L}{GEOMETRIE. BRIEFWECHSEL.}
\indent 4) dass die Unterscheidung in beiden Fällen
\begin{center}
\begin{tikzpicture}[x=.016cm,y=-.016cm,line width=.45pt]
\draw (64,2) circle[radius=3.3];
\draw (185,0) circle[radius=3.3];
\draw (0,102) circle[radius=3.3];
\draw (142,98) circle[radius=3.3];
\end{tikzpicture}\hspace{1.8cm}
\begin{tikzpicture}[x=.016cm,y=-.016cm,line width=.45pt]
\draw (0,0) circle[radius=3.3];
\draw (60,50) circle[radius=3.3];
\draw (194,76) circle[radius=3.3];
\draw (2,116) circle[radius=3.3];
\end{tikzpicture}
\end{center}
\noindent auf eine analytische Bedingung reducirt werde, so würde sich das auch thun lassen, obwohl ich die Entwicklung in diesem Augenblick nicht gegenwärtig habe .\enspace.\enspace.\enspace.\enspace.

\NumberedItem{6}
\begin{center}\textsc{Schumacher} an \textsc{Gauss}.\quad Altona, 19. April 1836.\end{center}
\HistoricCitation{[Briefwechsel zwischen \textsc{Gauss} und \textsc{Schumacher} III, Altona 1861, S. 39—40.]}
\indent \{Ich sehe aus Ihrer Antwort, mein theuerster Freund, mit Beschämung, dass ich mich sehr undeutlich ausgedrückt haben muss. Vor allem bitte ich Sie überzeugt zu seyn, dass es nie meine Absicht gewesen ist, dass Sie sich \textls[120]{abquälen} sollten, ein Missverständniss, das ich kaum möglich gehalten hätte, wenn ich auch vergessen habe hinzuzufügen, dass meine Bitten nur unter der Voraussetzung gemacht waren, dass Sie sie in einigen müssigen Augenblicken erfüllen konnten.

\indent Es kann jetzt nicht von dem, was ich wünschte, mehr die Rede seyn; es sei mir nur erlaubt historisch zu bemerken, dass ich in Bezug auf Ihren Brief vom 21. März (in dem Sie den Punct, der dem gesuchten Minimum bei einem Viereck \(abcd\) entspricht, in \textls[120]{einen} der 3 Durchschnittspuncte der Linien \(ab,\ cd;\ ac,\ bd;\ ad,\ bc\) legen) \textls[120]{diese 3 Puncte} aus den Gleichungen zwischen den Cosinussen und Sinussen von \(w\) abzuleiten suchte, aber nur, wenn alle inneren Winkel \(<180^{0}\), den Durchschnittspunct der Diagonalen finden konnte, und daher zu wissen wünschte, wie man diese 3 Puncte \textls[120]{durch Rechnung} aus den bekannten Gleichungen finden könnte.

\indent Ich wünschte ferner zu wissen, ob sich die gesuchten Coordinaten \(x,y\) aus den Gleichungen zwischen den Coordinaten finden lassen, wobei ich schon für das Viereck auf unübersehbare Complication stiess.

\indent Dass aus den Gleichungen für \(\sin\) und \(\cos\) der Winkel \(w\) mein Schluss mit Nothwendigkeit folge, hatte ich freilich nicht bewiesen, ich vermuthe aber
\clearpage\SourceLeaf{474}\ArithmeticHead{467}{R}{SCHUMACHER AN GAUSS.}
\noindent der Beweis wird sich aus dem Umstande, dass die Lage der Linie, von welcher die Winkel \(w\) gezählt werden, \textls[120]{beliebig} ist, herleiten lassen.

\indent Nachdem ich so das, was ich sagen wollte, erklärt habe, bitte ich Sie herzlich, meinen Wunsch als nicht gemacht anzusehen und weiter nicht zu beachten .\enspace.\enspace.\enspace.\enspace.\}

\NumberedItem{7}
\begin{center}\textsc{Gauss} an \textsc{Schumacher}.\quad Göttingen, 21. April 1836.\end{center}
\HistoricCitation{[Briefwechsel zwischen \textsc{Gauss} und \textsc{Schumacher} III, Altona 1861, S. 44—45.]}
\indent Es scheint, mein theuerster Freund, dass Sie den von mir ungeschickter Weise gebrauchten Ausdruck abquälen härter auslegen als er gemeint war. Ich übernehme gern (wie ich glaube öfters bewiesen zu haben) selbst weitläuftige Auseinandersetzungen, wenn ich weiss, Ihnen damit eine Gefälligkeit zu erweisen. Ich wollte nur sagen, dass ich vorher gewiss zu sein wünschte, \textls[120]{welches} der dunkle Punkt sei, um nicht zu riskiren, das zu thun, was die Römer hircum mulcere nannten.

\indent In der That bin ich es bei dem fraglichen Gegenstande noch immer nicht. Es lässt sich

\indent 1) aus Ihren Gleichungen \textls[120]{beweisen}, dass nur der Durchschnittspunkt zweier Seitenpaare ihnen Genüge leistet.

\indent Es scheint aber nach Ihrem Briefe, dass Sie auf die Entwicklung dieses Beweises, der die \textls[120]{Nothwendigkeit} zeigt, keinen Werth legen, und \textls[120]{in dieser Ungewissheit} enthalte ich mich jetzt der Entwicklung.

\indent Dann aber hat auch die analytische Auflösung, den Punkt zu finden, wo zwei durch zwei Paare gegebener Punkte gezogene Linien einander schneiden, gar keine Schwierigkeit. In der That, wenn es die beiden Paare \(a,b\) mit \(a',b'\) und \(a'',b''\) mit \(a''',b'''\) sind, so hat man die beiden Gleichungen
\[
\begin{aligned}
x(b'-b)+y(a-a')+a'b-ab'&=0\\
x(b'''-b'')+y(a''-a''')+a'''b''-a''b'''&=0,
\end{aligned}
\]
\noindent woraus \(x\) und \(y\) durch Elimination folgen.

\indent Und ebenso erhält man die beiden andern Durchschnittspunkte .\enspace.\enspace.\enspace.\enspace.
\par\vfill\noindent\hfill59\textsuperscript{*}
